% !TeX encoding = UTF-8
% !TeX program = pdflatex
% IEEE Conference LaTeX Template
% Based on IEEEtran.cls (conference mode)
% Compile: pdflatex → bibtex → pdflatex → pdflatex

\documentclass[conference]{IEEEtran}
\IEEEoverridecommandlockouts

\usepackage{cite}
\usepackage{amsmath,amssymb,amsfonts}
\usepackage{algorithmic}
\usepackage{graphicx}
\usepackage{textcomp}
\usepackage{xcolor}
\usepackage{booktabs}
\usepackage{hyperref}

\def\BibTeX{{\rm B\kern-.05em{\sc i\kern-.025em b}\kern-.08em
    T\kern-.1667em\lower.7ex\hbox{E}\kern-.125emX}}

\begin{document}

% ============================================================
% Title & Authors
% ============================================================
\title{Paper Title}

\author{
\IEEEauthorblockN{Author One\IEEEauthorrefmark{1}, Author Two\IEEEauthorrefmark{2}}
\IEEEauthorblockA{
    \IEEEauthorrefmark{1}Department, University, City, Country\\
    \IEEEauthorrefmark{2}Department, University, City, Country\\
    Email: \{author1, author2\}@example.com}
}

\maketitle

% ============================================================
% Abstract
% ============================================================
\begin{abstract}
Abstract content (150--250 words).
\end{abstract}

\begin{IEEEkeywords}
keyword1, keyword2, keyword3, keyword4, keyword5
\end{IEEEkeywords}

% ============================================================
% 1  Introduction
% ============================================================
\section{Introduction}

% Para 1: Task importance and real-world applications
% Para 2: Current progress — representative method families
% Para 3: Remaining gap — what existing methods fail to do and WHY
% Para 4: Proposed method — core mechanism and intuition
% Para 5: Contributions — prose lead-in + exactly 3 detailed bullets, each matching one innovation point

\section{Related Work}

% Theme-based structure, not paper-by-paper
\subsection{Theme A}
\subsection{Theme B}

% ============================================================
% 2  Method
% ============================================================
\section{Method}

% Opening paragraph (3-5 sentences): overall framework → what each of the
% three innovations solves → how they collaborate. No "Overview" subsection.
% Reference the architecture diagram: As shown in Fig.~X, ...

\subsection{<Innovation 1 Name>}

% Motivation → Design → Formula/Mechanism → Effect
% Optional: module diagram, algorithm pseudocode

\subsection{<Innovation 2 Name>}

% Motivation → Design → Formula/Mechanism → Effect

\subsection{<Innovation 3 Name>}

% Motivation → Design → Formula/Mechanism → Effect

% If the loss function is an independent innovation, give it a subsection.
% If it supports a specific module, include it within that module's subsection.

\subsection{Implementation Details}

% Optimizer, learning rate schedule, batch size, data augmentation, etc.

% ============================================================
% 3  Experiments
% ============================================================
\section{Experiments}

\subsection{Datasets and Metrics}

% Dataset size, classes, splits; metrics and direction (higher/lower is better)

\subsection{Experimental Setup}

% Framework, optimizer, LR schedule, batch size, epochs, input size,
% data augmentation, GPU, random seeds, baseline configuration

\subsection{Comparison with State-of-the-Art and Classical Methods}

% SOTA comparison: main results table + quantitative analysis
% Classical methods comparison

\subsection{Ablation Study}

% Numbered list 1), 2), 3)... analyzing each component / design choice / hyperparameter

\subsection{Visualization}

% Feature maps, attention maps, t-SNE, confusion matrices, error cases, etc.

% ============================================================
% 4  Conclusion
% ============================================================
\section{Conclusion}

% Paragraph 1: Full summary (problem → method → key results with numbers)
% Paragraph 2: Limitations + future work (specific limitations + actionable directions)

% ============================================================
% Acknowledgments (optional)
% ============================================================
% \section*{Acknowledgment}
% Acknowledgment text.

% ============================================================
% References
% ============================================================
\bibliographystyle{IEEEtran}
\bibliography{references}

\end{document}
